\documentclass[10pt,a4paper]{article}
\usepackage[margin=2.2cm]{geometry}
\usepackage{fontspec}
\usepackage{kotex}
\setmainfont{Apple SD Gothic Neo}
\setsansfont{Apple SD Gothic Neo}
\usepackage{amsmath,amssymb}
\usepackage{graphicx}
\usepackage{booktabs}
\usepackage{caption}
\usepackage[dvipsnames]{xcolor}
\usepackage[colorlinks=true,linkcolor=MidnightBlue,citecolor=MidnightBlue,
            urlcolor=MidnightBlue]{hyperref}
\usepackage{titlesec}
\titleformat{\section}{\large\bfseries}{\thesection}{0.6em}{}
\captionsetup{font=small,labelfont=bf}
\setlength{\parskip}{0.35em}
\setlength{\parindent}{0pt}

\usepackage{mdframed}
\newmdenv[linewidth=0.4pt,linecolor=gray!60,backgroundcolor=gray!5,
          innertopmargin=6pt,innerbottommargin=6pt,skipabove=8pt,skipbelow=8pt]{claimbox}
\newcommand{\claim}[2]{\begin{claimbox}\textbf{주장.} #1\par\smallskip
  \textcolor{gray!70}{\textbf{반증 조건.} #2}\end{claimbox}}

\title{\bfseries 오염되지 않은 지도\\[2pt]
\large 합동승은 같은-IP 조회가 아니다 --- 그리고 겹침행은 반대 방향을 가리켰다}
\author{월드모델 실험실 · 노트 $825$}
\date{2026-08-07}
\begin{document}\maketitle

\section{한 줄}

티처 세션의 최우선 지적을 이행했다 --- \emph{``$822$(b)의 합동승 도메인이
정확히 겹침 큰 쪽이다. 같은-IP 교차 노출이 가장 유력한 제$3$ 변수이고 지금
검사 가능하다. 통과 전엔 고객 덱에 넣지 마라.''} 세계애니 유보 $300$ 을
겹침(만화 학습과 같은 제목 $86$행)/비겹침($214$행)으로 갈라 합동$-$릿지 차를
재채점했다.

\begin{center}
\begin{tabular}{lrrrr}
\toprule
세계애니 & $n$ & 합동 & 릿지 & $\Delta$ \\
\midrule
전체($822$) & $300$ & $0.5935$ & $0.4508$ & $+0.1427$ \\
\textbf{비겹침} & $214$ & $0.5751$ & $0.4288$ & $\mathbf{+0.1464}$ ($2\sigma$ 의 $4.1$배) \\
겹침(만화 원작) & $86$ & $0.5972$ & $0.5383$ & $+0.0590$ \\
\bottomrule
\end{tabular}
\end{center}

\textbf{판정: 통과 --- 오염 가설 기각.} 합동의 우위는 같은-IP 노출 없이
성립한다. 게임(겹침 $4$행) · 애니($12$행 · 한글$\leftrightarrow$로마자 하한)는
검정 불능으로 각주 병기.

\begin{figure}[t]\centering
\includegraphics[width=0.97\textwidth]{figs/map.pdf}
\caption{왼쪽: 세 부분집합의 합동 대 릿지. 오른쪽: 가설 재판의 요약.}
\end{figure}

\section{부 예측이 정반대로 틀렸고, 그것이 더 배웠다}

\claim{\textbf{오염 기제라면 겹침행에서 $\Delta$ 가 커야 했다}(내 부 예측
$60\%$). 실측은 \textbf{절반}($+0.059$)이다 --- 같은-IP 행(만화 원작이 있는
애니화)은 자기-도메인 정보만으로도 맞힐 수 있는 쉬운 행이라 \textbf{릿지가
따라잡는 곳}이었다(릿지 $0.5383$ 대 비겹침 $0.4288$). 깊이 $12$ 트리의 퍼지
IP 조회 가설과 반대 방향. 따라서 $822$ 의 `합동 수혜 도메인의 지도'는
살아 있던 유일한 오염 가설을 통과했고 \textbf{판권/고객 대화에 들고 갈 자격}을
얻었다 --- 표본 공유 실증(합동승 $5$)도 유지된다.}
{애니는 언어-정합 매칭(로마자$\leftrightarrow$한글)이 생겨 겹침행 $\geq 20$
이 확보되면 $1$회 재검한다(결정 규칙). 판 $\rho$ $0.4689$ 그대로.}

\end{document}
